% ============================================================
% 第 5 章 三相变压器（教材第二篇，印刷页 48--59）
% 节名：第一节 三相变压器组和三相心式变压器；第二节 三相变压器的联结组；
%       第三节 三相变压器中励磁电流、主磁通和感应电动势波形；第四节 标幺值
% ============================================================
\chapter{三相变压器}

由于电能的产生、传输和分配绝大部分采用三相制，三相变压器的应用十分广泛。实际运行时三相变压器的电压和电流基本上是对称的，对称三相电压或电流大小相等、相位互差 $120^\circ$，因此可取三相中的一相来分析计算，即将三相问题简化成单相问题。这样，第 4 章的基本方程、等效电路和相量图对三相变压器仍然适用，本章只介绍三相变压器的几个特殊问题：磁路结构、联结组、波形与标幺值。

\section{三相变压器组和三相心式变压器}

三相变压器按铁芯结构分类，通常有三相变压器组和三相心式变压器两种。

\subsection{三相变压器组}
三相变压器组由三个相同的单相变压器按一定的连接方式组合而成。三相变压器组的\textbf{三相磁路各自独立}。当一次绕组外加三相对称电压时，三相磁路中的主磁通和三相绕组中的励磁电流基波都是三相对称的。

\subsection{三相心式变压器}
如果把三个单相变压器的铁芯合并，每相的高、低压绕组放在一起，并将三相绕组分别套装在外侧的三个铁芯柱上，则磁路仍保持三相对称。当三相绕组外接对称三相电压时，三相磁通亦为三相对称，且三相磁通之和为零：
\begin{equation}
\dot\Phi_A+\dot\Phi_B+\dot\Phi_C=0 .
\end{equation}
也就是说，中间的铁芯柱中没有磁通经过，可以将它省去；如果再使三个铁芯柱布置在同一平面上，便得到三相心式变压器的铁芯。

三相心式变压器的磁路是\textbf{相互联系}的，相当于电路中的 Y 形连接，任一相的磁通都要借助另外两相的磁路闭合。由于三相磁路的长度不同（中间芯柱的磁路较短、磁阻较小），三相磁路的磁阻不同：在三相对称电压作用下，三相磁路中的主磁通相同，但产生主磁通所需的励磁电流却不同，中间的励磁电流稍小一些。工程上常取它们的平均值作为实际的励磁电流。由于磁路的不对称程度很小，三相励磁电流差别不大，对运行性能影响也很小，通常可不予考虑。

\begin{center}
\begin{tabular}{@{}p{2.3cm}p{3.5cm}p{4.0cm}p{3.6cm}@{}}
\toprule
项目 & 三相变压器组 & 三相心式变压器 & 说明 \\
\midrule
磁路 & 三相独立 & 三相相互联系（Y 形） & 决定 3 次谐波磁通能否流通 \\
磁阻 & 三相等 & 三相不等（中柱小） & 励磁电流略有差异，取平均值 \\
硅钢片用量 & 多 & 少（省去中间柱） & 心式质量轻、造价低 \\
制造与运输 & 便于制造、运输，可降低备用容量 & 大容量时运输受限 & 组式多用于巨型大容量 \\
\bottomrule
\end{tabular}
\end{center}

心式变压器的优点是：硅钢片用量少，质量轻，造价便宜，应用比较广泛。三相变压器组便于制造和运输，也可使用户的备用容量降低，故多用于巨型大容量的变压器。

\section{三相变压器的联结组}

\subsection{绕组标记和三相绕组的连接方法}
变压器高、低压绕组出线端的标记，国家标准有统一规定。单相变压器的高压绕组出线端以大写字母 $A$、$X$ 标记，低压绕组出线端以小写字母 $a$、$x$ 标记，其中 $A$、$a$ 表示首端，$X$、$x$ 表示末端。三相变压器的高压绕组首端标以 $A$、$B$、$C$，末端标以 $X$、$Y$、$Z$；低压绕组的首端标以 $a$、$b$、$c$，末端标以 $x$、$y$、$z$。

三相绕组连接方法很多，最常用的有两种：
\begin{itemize}
  \item \textbf{星形接法}：把三相绕组的末端连在一起，把它们的首端引出来；
  \item \textbf{三角形接法}：又分两种——按 $ax$—$by$—$cz$ 顺序连接称为\textbf{顺序三角形接法}；按 $ax$—$cz$—$by$ 顺序连接称为\textbf{逆序三角形接法}。
\end{itemize}

GB 1094.1—1996《电力变压器 第 1 部分：总则》规定：星形接法高压绕组用大写 $Y$ 表示、低压绕组用小写 $y$ 表示，有中性线引出时用 $YN$ 或 $yn$ 表示；三角形接法高压绕组用 $D$ 表示、低压绕组用 $d$ 表示。\textbf{高压侧接法写在前面，低压侧接法写在后面}。将高、低压绕组作不同接法组合，可得到变压器的四种基本联结组：$Yy$、$Yd$、$Dy$、$Dd$，其中 $Y(y)$ 联结组也包括有中性线引出的 $YN(yn)$ 联结组。前两种是国产变压器最常用的联结组，最后一种在国产变压器中基本不用。

\subsection{单相变压器高、低压绕组电动势之间的相位关系}
单相变压器中，高、低压绕组的电动势要么同相、要么反相，取决于两个绕组的绕制方向和标记方法。由于高、低压绕组被同一磁通 $\Phi$ 交链，其感应电动势有一定的极性关系：当高压绕组某一端瞬时电位为正时，低压绕组也有一电位为正的对应端，这两个对应的同极性端点称为\textbf{同名端}，用符号"$\cdot$"表示。

高、低压绕组的绕制可以同向也可以反向，任一出线端可以作为首端也可以作为末端。设绕组电动势正方向均规定为从首端指向末端，则根据绕制方向和标记方法有四种组合：
\begin{itemize}
  \item 绕向相同、标记相同：同名端在相同端，$\dot E_A$ 与 $\dot E_a$ 同相；
  \item 绕向相同、标记相反：同名端在不同端，$\dot E_A$ 与 $\dot E_a$ 反相；
  \item 绕向相反、标记相同：同名端在不同端，反相；
  \item 绕向相反、标记相反：同名端在相同端，同相。
\end{itemize}

规则可概括为：\textbf{若首—首（或末—末）为同名端，则 $\dot E_A$ 与 $\dot E_a$ 同相；若首—末（或末—首）为同名端，则反相}。

变压器中通常用\textbf{时钟法}表示高、低压绕组电动势的相位关系：高压绕组电动势大，代表时钟的\textbf{长针}，并认为它始终指向 0 点钟位置；低压绕组电动势小，代表\textbf{短针}。同相时相当于 0 点钟，反相时相当于 6 点钟。因此单相变压器只有两种情况，分别用 $Ii0$ 和 $Ii6$ 表示（$I$ 或 $i$ 表示单相变压器，数字表示相位关系）。国家标准规定 $Ii0$ 为标准联结组。

\subsection{三相变压器高、低压绕组线电动势之间的相位关系}
三相变压器除相电动势的相位关系外，还有线电动势的相位关系。后者除决定于绕向和标记方法外，还决定于\textbf{三相绕组的连接方法}。

\subsubsection{Yy 联结组}
若高压绕组 $A$、$B$、$C$ 相首端分别与低压绕组 $a$、$b$、$c$ 相首端为同名端，则高压绕组的相电动势 $\dot E_A$、$\dot E_B$、$\dot E_C$ 分别与低压绕组相对应的相电动势 $\dot E_a$、$\dot E_b$、$\dot E_c$ 同相位。由于两侧都是 Y 接，相对应的线电动势也同相。

分析线电动势时用 $\dot E_{BC}=\dot E_B-\dot E_C$，方向由 $\dot E_C$ 的箭头指向 $\dot E_B$ 的箭头；同理 $\dot E_{bc}=\dot E_b-\dot E_c$。两电动势相量平行，夹角为零，具有相同相位。用时钟法表示时，在三相系统中高压绕组的线电动势表示长针并指向 0（或 12），低压绕组的线电动势表示短针，故相位关系相当于 0 点钟，该联结组别为 \textbf{Yy0}。

如果将高压绕组 $A$、$B$、$C$ 相首分别与低压绕组的 $a$、$b$、$c$ 相\textbf{尾}为同名端，则联结组别变成 \textbf{Yy6}。

\subsubsection{Yd 联结组}
若高压绕组的相电动势 $\dot E_A$、$\dot E_B$、$\dot E_C$ 分别与低压绕组相对应的相电动势同相位，但由于高压绕组为 Y 接，\textbf{线电动势与相电动势之间有 $30^\circ$ 的相位移}；低压绕组为 d 接，线电动势等于相电动势。因此虽然相电动势相位相同，但由于接法不同，两侧相对应线电动势的相位不会相同。

以 $\dot E_{BC}$ 与 $\dot E_{bc}$ 为例：高压侧为 Y 接，$\dot E_{BC}$ 与图 5-5 相同；低压侧由 $\dot E_{bc}=\dot E_b-\dot E_c$，其相量与 $\dot E_b$ 相反，超前 $30^\circ$，相当于超前一个钟点。按时钟表示法，相位关系为 11 点钟，故该联结组别为 \textbf{Yd11}。

如果将同名端改为非同名端，则联结组别变成 \textbf{Yd5}；如果其他不变、仅将低压绕组由逆接改为顺接，则变成 \textbf{Yd1}。

\subsubsection{12 个组别}
三相变压器的联结组共有 12 个，相邻两个联结组相位差 $30^\circ$，12 个联结组构成圆的一周，和时钟的 12 个整点相同——这实质上是联结组用时钟法表示的原因。

在 Yy0（或 Yy6）基础上，若高压绕组三相标号 $AX$、$BY$、$CZ$ 不变，把低压绕组三相标号 $ax$、$by$、$cz$ 按顺时针方向旋转 $120^\circ$（$ax$ 被 $by$ 取代、$by$ 被 $cz$ 取代、$cz$ 被 $ax$ 取代），相当于短针顺时针转动 4 个钟点，得到 Yy4（或 Yy10）；逆时针旋转 $120^\circ$ 则得到 Yy8（或 Yy2）。故 \textbf{Yy 联结组具有 0、4、8、6、10、2 这 6 个偶数组别}；同理 \textbf{Yd 联结组具有 11、3、7、5、9、1 这 6 个奇数组别}。

Yy 联结组的某一组号通常只有一种连接方式；而 Yd（或 Dy）联结组的某一组别通常有\textbf{顺接和逆接两种方式}。为制造和使用方便，国家标准规定我国只生产 \textbf{Yyn0、Yd11、YNd11、YNy0 以及 Yy0} 共五种标准联结组，其中前三种最常用。

\begin{center}
\begin{tabular}{@{}lll@{}}
\toprule
联结组 & 钟点数 & 特点与用途 \\
\midrule
Yyn0 & 0 & 配电变压器最常用，可引出中性线供三相四线制 \\
Yd11 & 11 & 大容量输电、降压变压器常用（一侧三角形可抑制 3 次谐波） \\
YNd11 & 11 & 高压侧中性点可接地，用于系统联络 \\
YNy0 & 0 & 高压侧中性点接地 \\
Yy0 & 0 & 仅中小容量 \\
\bottomrule
\end{tabular}
\end{center}

\corepoint{判断联结组的通用步骤（考场上照做即可）：① 按同名端标出一次侧三相相电动势相量（互差 $120^\circ$）；② 按同名端定出二次侧相电动势与一次侧相电动势的同相/反相关系；③ 按接法（Y 或 D）由相电动势合成线电动势；④ 比较同一对线电动势（如 $E_{AB}$ 与 $E_{ab}$）的夹角，除以 $30^\circ$ 即为钟点数。\textbf{不要背结论，把四步画一遍}——组别题一旦背结论必错。}

\section{三相变压器中励磁电流、主磁通和感应电动势波形}

由第 1 章第五节：磁路饱和时，若磁通为正弦波，则励磁电流为尖顶波；尖顶波的主要成分是基波与 3 次谐波。3 次谐波电流频率为基波 3 倍，在三相系统中\textbf{三相同相位}：
\begin{equation}
i_{m3A}=I_{m3}\sin3\omega t,\quad
i_{m3B}=I_{m3}\sin3\left(\omega t-120^\circ\right)=I_{m3}\sin3\omega t,\quad
i_{m3C}=I_{m3}\sin3\omega t .
\end{equation}
反之，当励磁电流为正弦波时，铁芯中的磁通为平顶波，其主要成分同样是基波与 3 次谐波，且 3 次谐波磁通也同相位。

对单相变压器，3 次谐波电流可以流通，因此在正弦电源电压作用下主磁通及感应电动势都接近正弦波。但在三相变压器中，3 次谐波电流和 3 次谐波磁通能否流通，分别要由\textbf{绕组的联结方式}和\textbf{磁路结构}决定。

\subsection{YNy 联结组}
高压绕组有中性线时，中性线上的电流为 $i_N=i_{mA}+i_{mB}+i_{mC}$。由于励磁电流中三相基波分量之和为零（基波电流仅在三相绕组内流通，不经过中性线），所以中性线上的电流只有 3 次谐波电流：
\begin{equation}
i_N=i_{m3A}+i_{m3B}+i_{m3C}=3I_{m3}\sin3\omega t .
\end{equation}
即励磁电流中的 3 次谐波分量在三相绕组与中性线之间形成闭合回路，可以在绕组中流通。此时情况与单相变压器相似：外接正弦三相电压时，即使磁路饱和，铁芯中的磁通波形和绕组中的感应电动势波形仍接近正弦波。由于中性线上的谐波电流会产生损耗和通信干扰，YNy 联结组很少采用。

\subsection{Yy 联结组}
去掉中性线后，励磁电流中的 3 次谐波分量不能流通，励磁电流波形接近正弦波；于是主磁通被迫畸变，成为\textbf{平顶波}。3 次谐波磁通能否流通由磁路结构决定。

\subsubsection{三相变压器组}
三相磁路分开、各自独立，3 次谐波磁通可以和基波磁通一样在铁芯内闭合。基波磁通产生基波电动势 $e_1$，3 次谐波磁通产生 3 次谐波电动势 $e_3$，两者的相位都滞后于各自的磁通 $90^\circ$。虽然 $\Phi_{m3}<\Phi_{m1}$，但 $e_3\propto3f\Phi_{m3}$，所以 3 次谐波电动势相当大，可达基波电动势的 45\%～60\%。因此 $e_1$ 与 $e_3$ 叠加后，每相绕组的合成电动势发生较大畸变，成为\textbf{尖顶波}；尖顶波使相电动势出现相当大的尖峰值，对绕组绝缘不利。因此\textbf{三相变压器组不宜采用 Yy 联结组}。在三相线电动势中，3 次谐波电动势相互抵消，故仍为正弦波。

\subsubsection{三相心式变压器}
铁芯磁路相互联系，相当于星形接法的磁路，3 次谐波磁通在铁芯磁路中受到抑制，不得不从空气（或油）、油箱壁等\textbf{漏磁路}中闭合。漏磁路磁阻较大，$\Phi_{m3}$ 很小，因而 $e_3$ 也很小，主磁通和感应电动势仍接近正弦波。但 3 次谐波磁通流经漏磁路的油箱壁等铁磁部件，会引起附加涡流损耗、局部发热并降低效率，所以采用 Yy 联结组的三相心式变压器容量不宜太大，通常不超过 1800\,kV·A。

\subsection{Dy 和 Yd 联结组}
\begin{itemize}
  \item \textbf{Dy}：一次绕组为三角形连接，3 次谐波电流可以流通，故在正弦三相电压下励磁电流为尖顶波，主磁通及 $e_1$、$e_2$ 皆为正弦波，与单相变压器相似。
  \item \textbf{Yd}：一次绕组为星形连接，3 次谐波电流不能流通，一次绕组中励磁电流接近正弦波，故铁芯磁路中的主磁通含有 3 次谐波分量；该分量在二次绕组中感应 3 次谐波相电动势 $e_{23}$，而二次绕组为三角形连接，三相的 $e_{23}$ 将在闭合三角形绕组中产生 3 次谐波\textbf{环流} $i_{23}$。
\end{itemize}
由于一次绕组为星形连接、3 次谐波电流不能流通，不可能产生 3 次谐波磁动势与 $N_2\dot I_{23}$ 相平衡，故二次绕组中的环流就起到了励磁电流的作用。磁路中的主磁通由与它相交链的\textbf{合成磁动势}产生，而合成磁动势不必非由一个绕组提供；产生正弦主磁通所必须的基波励磁电流与 3 次谐波电流由一、二次绕组共同提供，与单独由一次绕组提供效果相同。因此 Yd 与 Dy 效果一样，主磁通和每相绕组感应电动势仍接近正弦波形。

三角形联结组中的 3 次谐波环流甚小，对运行影响不大，且能改善主磁通和感应电动势波形，因此较大容量的变压器常将其中一侧绕组接成三角形。

\begin{center}
\begin{tabular}{@{}p{1.5cm}p{2.5cm}p{3.4cm}p{4.2cm}@{}}
\toprule
联结组 & 3 次谐波电流 & 3 次谐波磁通 & 相电动势波形 \\
\midrule
YNy & 经中性线流通 & 很小 & 接近正弦 \\
Yy（组式） & 不能流通 & 可在铁芯内闭合 & 尖顶波（$e_3$ 达 45\%～60\%），不宜采用 \\
Yy（心式） & 不能流通 & 受抑制，经漏磁路闭合 & 接近正弦，但油箱附加损耗，容量 $\le1800$\,kV·A \\
Dy & 一次 D 内流通 & 很小 & 正弦 \\
Yd & 二次 D 内环流 & 由两侧合成磁动势建立 & 正弦，大容量常用 \\
\bottomrule
\end{tabular}
\end{center}

\corepoint{这一节的判断链只有三步：① 有没有三角形接法或中性线 $\Rightarrow$ 3 次谐波电流能否流通；② 磁路是组式（独立）还是心式（联系）$\Rightarrow$ 3 次谐波磁通能否在铁芯内闭合；③ 3 次谐波磁通大不大 $\Rightarrow$ 相电动势是否尖顶。记住"线电动势中 3 次谐波互相抵消、永远正弦"，这是常考的一句话陷阱。}

\section{标幺值}

变压器和电机中，电压、电流、阻抗和功率等物理量可用实际值，也常用\textbf{标幺值}表示与计算。标幺值也称相对值，是\textbf{同单位}的实际物理量与其基值之比：
\begin{equation}
\text{标幺值}=\frac{\text{实际值}}{\text{基值}} .
\end{equation}
本书中基值用下角标"b"表示，标幺值用右上角标"*"表示。标幺值乘以 100\% 便得百分值。

\subsection{基值}
应用标幺值时先要选定基值。变压器和电机的工程计算中，通常选定额定值作为基值。

对单相系统：
\begin{equation}
S_b=P_b=Q_b=U_bI_b=U_NI_N=S_N,\qquad
Z_b=R_b=X_b=\frac{U_N}{I_N}=\frac{U_N^{2}}{S_N}.
\end{equation}
对三相系统：
\begin{equation}
U_{bL}=U_{NL}=U_N,\qquad
S_b=P_b=Q_b=\sqrt3U_{NL}I_{NL}=S_N,\qquad
Z_b=\frac{U_{N\phi}}{I_{N\phi}}=\frac{U_{NL}^{2}}{S_N},
\end{equation}
其中下标 $\phi$ 表示相值、$L$ 表示线值；$Z_b$ 在星形接法时为 $U_{NL}^{2}/S_N$，三角形接法时相应取相值形式。

将各物理量的实际值与对应基值之比，即得各物理量的标幺值。注意：\textbf{一次绕组和二次绕组的基值分别选取各自的基值}。例如单相变压器一次绕组电压和电流的标幺值为 $U_1^{*}=U_1/U_{1N}$、$I_1^{*}=I_1/I_{1N}$；三相变压器一次绕组线电压的标幺值为 $U_{1L}^{*}=U_{1L}/U_{1N}$。

\subsection{标幺值的特点}

\subsubsection{额定电压、电流和容量的标幺值为 1}
由于电压、电流和容量的基值都选为额定值，故
\begin{equation}
U_N^{*}=I_N^{*}=S_N^{*}=1 .
\end{equation}

\subsubsection{在三相系统中，相、线标幺值相等}
无论电压还是电流，其相、线标幺值相等。以星形接法为例，
\begin{equation}
U_L^{*}=\frac{U_L}{U_N}=\frac{\sqrt3U_\phi}{\sqrt3U_{N\phi}}=\frac{U_\phi}{U_{N\phi}}=U_\phi^{*},\qquad
I_L^{*}=I_\phi^{*}.
\end{equation}
这一条极其实用：\textbf{用标幺值计算时不必区分相值与线值}，也就不用反复乘除 $\sqrt3$。

\subsubsection{单、三相功率计算公式相同}
用标幺值表示时，单相与三相系统的视在功率、有功功率和无功功率表达式相同，且不存在"线"的概念：
\begin{equation}
S^{*}=U^{*}I^{*},\qquad P^{*}=U^{*}I^{*}\cos\varphi,\qquad Q^{*}=U^{*}I^{*}\sin\varphi,
\end{equation}
式中 $\varphi$ 为功率因数角。

\subsubsection{归算前后的值相等}
以二次绕组电压和阻抗为例：
\begin{equation}
U_2'^{*}=\frac{U_2'}{U_{1N}}=\frac{kU_2}{kU_{2N}}=\frac{U_2}{U_{2N}}=U_2^{*},\qquad
Z_2'^{*}=\frac{k^{2}Z_2}{k^{2}Z_{2b}}=Z_2^{*}.
\end{equation}
即\textbf{用标幺值表示时，归算前后的值相等}——这意味着用标幺值计算可以\textbf{不必再进行归算}，这是标幺值最大的便利。

\subsubsection{变压器基本方程组的表达形式不变}
用标幺值表示，变压器基本方程组的表达形式保持不变，并且归算值上角标的"$'$"可以保留也可以去掉。

\subsubsection{便于比较不同容量的电机}
对于不同容量的变压器或电机，用标幺值表示时各参数和典型性能数据通常都在一定范围内，便于比较和分析。例如电力变压器漏阻抗的标幺值约为 4\%～8\%，空载电流的标幺值约为 1\%～5\%。

\corepoint{标幺值最容易错的三处：① 基值必须"成对"选取——$U_b$ 与 $I_b$ 要满足 $U_bI_b=S_b$，只改一个就会使阻抗基值错；② 一次侧、二次侧各有自己的基值，$U_{2b}=U_{2N}$（不是 $kU_{1N}$ 的换算值），这样才能自动消去 $k$；③ 三相系统选\textbf{线值}作基值时，$Z_b=U_{NL}^2/S_N$；若选相值作基值，$Z_b=U_{N\phi}^2/(S_N/3)=U_{N\phi}/I_{N\phi}$，两者结果相同，但不能混着用。}
